\chapter{Conclusion and Future Work}\label{ch:conclusion}

\section{Conclusion}
This report tested two ideas for cheaper training of a small neural network under one strict protocol: counted FLOPs to a fixed validation accuracy, equal tuning of every arm, paired seeds, and hypotheses locked before the confirmatory runs.

The first idea, an auxiliary meeting loss derived from simple difference target propagation, did not save compute. It was refuted by a pre-registered test ($+36.9\%$ and $+35.5\%$ more FLOPs than a tuned baseline), and a set of controls, variants and mechanism checks showed that the result is not an artefact of one detail: the signal is well aligned with the gradient, giving it more strength did not help, giving it an exact inverse next to backpropagation did not help, and with backpropagation cut the learned inverse did no better than a random one, and in ReLU the replacement trained only with a bounded signal, at more steps and no saving. The evidence is consistent with a redundancy explanation.

The second idea, back-propagating only through the highest-loss rows of each batch, did save compute. In a second pre-registered test it needed $35.1\%$ and $38.4\%$ fewer FLOPs than the tuned baseline and $30.2\%$ and $35.5\%$ fewer than a tuned small-batch baseline, in all 20 seeds. In a third, a leaner variant (keep $20\%$, no warm-up) needed $48.5\%$ and $48.8\%$ fewer FLOPs than the baseline and $46.5\%$ and $53.6\%$ fewer than a tuned Selective-Backprop reimplementation, whose cost is dominated by extra selection passes.

The protocol itself produced two lessons: tuning on two seeds is unreliable (two tunings of the same baseline tied at selection and differed by up to $34\%$ on fresh seeds), and counted FLOPs, optimiser updates and time per step can rank methods differently. Everything is limited to one MLP on MNIST and to counted FLOPs.

\section{Future work}\label{sec:future}
In the order in which they would most increase the weight of the results:
\begin{enumerate}
\item \textbf{A second setting}: a second dataset (Fashion-MNIST) and a small convolutional network, each with its own tuned baseline and a new pre-registered plan, repeating the three confirmatory comparisons. This requires a selective backward pass for layers other than plain linear ones (for example through autograd hooks on the kept rows).
\item \textbf{Wall-clock time to target} for A1, S30, S20w0 and Selective-Backprop on CPU and GPU, several widths and thread counts, with the full time-to-target measured and not derived. Without it, ``saves compute'' means FLOPs only.
\item \textbf{A completed related-work check}: the literature review of \cref{ch:background} covers the closest methods, but a systematic search (including for earlier selective-backward and top-$k$-loss variants, and for auxiliary target-propagation signals) should be completed and the bibliography re-checked before any submission (\cref{app:reproducibility}; it was verified entry by entry on 25 November 2026).
\item \textbf{The remaining mechanism checks for the meeting loss}: a deeper network with SGD and momentum (check C); a replacement test in a network whose upper half dominates the cost, the only setting in which the cut could pay (here the upper half is $13\%$ of the forward cost); and one shared, learning-rate-tuned $\tanh$ baseline for all $\tanh$ comparisons.
\item \textbf{The effect of removing the warm-up}: a mechanism is not known; optimiser-state and early-ranking explanations should be separated by dedicated experiments.
\item \textbf{Selective-Backprop}: search over $\beta$, history length and priming; an implementation that caches the selection forward pass; a cost model with time and not only FLOPs.
\item \textbf{More reliable tuning}: re-evaluating the top trials on fresh seeds before selection, and a larger number of tunings for every arm and not only for the baseline.
\end{enumerate}

\paragraph{From report to paper.} The material divides naturally into a focused paper, with the selective backward pass (Studies~2 and 3) as its main contribution and the meeting-loss negative result as a short motivating section, and the protocol and its lessons as a methodological contribution. The paper would need items 1--3 above. A workshop version could be written with the current evidence if the limitations of \cref{sec:limitations} are stated prominently.
